% ---------------------------------------------------------------------
% Section 2
% ---------------------------------------------------------------------

\begin{cnread}
本节把「从理想点 $q$ 到普通点 $p$ 的一族穿线路径」严格写下来，并回答一个
具体问题：既然这些路径的起点终点都一样，它们在费曼图、\ACS{}、切触空间三个
视角下分别是什么？结论是：它们恰好是算子求值映射的一条纤维
$\rho^{-1}(g_{p_0\to p})$；在 \ACS{} 里摊成一条水平线段（裸荷随 $M\to\infty$
指数发散），在切触空间里则落在同一根 Reeb 纤维的\textbf{不同高度}上。
本节只用 Paper~0 的第 7--9 章，不需要任何物理字典。
\end{cnread}

\subsection{Two kinds of ideal point}

Let $\E_0=(\HH^2,g_{\mu,\lambda},a)$ be the basic model with
$a(x,y)=-x/y$.  In horocyclic coordinates $u=x$ and $y=e^{\lambda v}$ one has
$a=-ue^{-\lambda v}$, and the two affine motions of Paper~0,
Proposition~3.6, read
\begin{equation}\label{eq:grid-motions}
  \mathsf X_s(x,y)=(x-sy,\,y),
  \qquad
  \mathsf Y_k(x,y)=\Bigl(x,\frac yk\Bigr),
  \qquad
  a\circ\mathsf X_s=a+s,\quad a\circ\mathsf Y_k=ka .
\end{equation}
The ideal boundary $\pt_\infty\HH^2=\R\cup\{\infty\}$ carries two kinds of point,
and the distinction is used constantly below.

\begin{enumerate}[label=(\roman*)]
  \item The point at infinity, $q=i\infty$, i.e. $y\to+\infty$ or
  $M=\lambda v\to-\infty$.  This is the common base point of the
  Borel subgroup $B_\infty=\Stab_{\PGL_2}(\infty)$: every affine motion fixes
  it, so it is the natural origin of an affine history, and in the physical
  reading it is the ultraviolet, asymptotically free end.
  \item A finite ideal point $q=x_\infty\in\R$, i.e. $y\to0^{+}$ or
  $M\to+\infty$.  This is the infrared end, reached by unbounded positive
  scaling $\mathsf Y_{e^{M}}$ with $M\to+\infty$.
\end{enumerate}

\begin{remark}[Ultraviolet versus infrared end]
Whether one marches from $q=i\infty$ inward to $p$, or reads the same word
backwards from $p$ towards a finite ideal point, is a choice of direction, not a
choice of geometry: the two readings are the two ends of one word.  Paper~0
keeps both ends in view because the multiplicative charge is the only charge
that descends to the motion group (Theorem~8.6), so only $M$ can be compared
across histories without an order convention.
\end{remark}

\subsection{The family of histories is one operator fiber}

\begin{remark}[The ideal point is a direction, not a starting point]
\label{rem:q-is-a-direction}
$B_\infty$ fixes $q=i\infty$, so no affine motion carries $q$ to an interior
point of $\HH^2$: a history cannot literally start at the ideal point.  What the
phrase ``a history from $q$ to $p$'' means in this note, and what it means in
Paper~0's grid, is a history that starts at an ordinary point $p_0$ deep in the
ultraviolet region --- large $y_0$, near $q$ --- and ends at $p$.  The
formulation below therefore fixes $p_0$ explicitly.  Its distance from $q$ enters
only through the total multiplicative charge, and every invariant used in this
note is a function of the charges, not of how deep the start was placed.
\end{remark}

Fix an ordinary starting point $p_0=(x_0,y_0)\in\HH^2$ and an ordinary terminal
point $p=(x_p,y_p)\in\HH^2$.  A \emph{comb history} from $p_0$ to $p$ is a finite
chronological word
\begin{equation}\label{eq:comb-history}
  \gamma=(\cM_{m_1},\cA_{t_1},\cM_{m_2},\cA_{t_2},\dots,\cM_{m_n},\cA_{t_n}),
  \qquad m_i,t_i\in\R ,
\end{equation}
read from left to right, whose composite motion carries $p_0$ to $p$.  Here
$\cM_{m}$ denotes scaling by $e^{m}$ and $\cA_t$ translation by $t$; the word is
the left-expanded comb of Paper~0, Subsection~2.1.  In the comb notation of this
section we set the multiplicative step constant to $\lambda=1$, so that the
log-scale parameter $v$ and the multiplicative charge $M=\lambda v$ coincide;
the general case is recovered by the substitution $M=\lambda v$, and every
formula below is homogeneous in the charges.

\begin{proposition}[The $p_0\to p$ family is one operator fiber]
\label{prop:operator-fiber}
Let $g_{p_0\to p}\in B_\infty$ be the affine motion carrying $p_0$ to $p$.  Then
$\gamma$ in \eqref{eq:comb-history} carries $p_0$ to $p$ if and only if its
composite is $g_{p_0\to p}$.  Consequently the histories joining $p_0$ to $p$
are exactly one fiber of the evaluation map
$\rho:\Hist\to\Aff^{+}(1,\R)$ of Paper~0, Section~7.
\end{proposition}

\begin{proof}
$B_\infty$ acts simply transitively on $\HH^2$: it acts transitively because
$\mathsf Y_k$ and $\mathsf X_s$ of \eqref{eq:grid-motions} already carry any
point to any other, and freely because a nonidentity affine map has at most one
fixed point in $\HH^2$ (Paper~0, Proposition~5.4 with $b=0$; the pure
translation has none).  Hence exactly one affine motion carries $p_0$ to $p$,
and a word realizes that motion exactly when it equals $g_{p_0\to p}$ in the
group.
\end{proof}

Writing $g_{p_0\to p}(z)=e^{M_*}z+\xi_*$, the two invariants of the fiber are the
total multiplicative charge $M_*=\lambda\sum_i m_i$ and the total target
translation $\xi_*$.  In terms of the assignment,
\begin{equation}\label{eq:fiber-invariants}
  e^{M_*} = \frac{y_p}{y_0},
  \qquad
  a(p)=e^{M_*}a(p_0)+\xi_* = -\frac{x_p}{y_p}.
\end{equation}

\subsection{Comb diagrams and charge accounting}

The affine cocycle of Paper~0, Proposition~7.1, expresses the terminal
translation of a comb in two equivalent ways,
\begin{equation}\label{eq:comb-cocycle}
  \xi_\gamma=\sum_{i=1}^{n}t_i\,e^{\,M_{i+1:n}}
  = e^{M_*}\sum_{i=1}^{n}t_i\,e^{-M_i},
  \qquad M_{i:n}=\sum_{j=i}^{n}\lambda m_j .
\end{equation}
Read as a diagram, \eqref{eq:comb-cocycle} is a time-ordered Feynman diagram
with a single fermion line.

\begin{enumerate}[label=(\alph*)]
  \item The vertical segments $\cM_{m_i}$ are free propagation between
  vertices: they advance the scale by $e^{m_i}$ and cost no additive charge.
  \item The horizontal segments $\cA_{t_i}$ are vertex insertions at the scale
  shell $M_i$: each emits or absorbs one quantum and adds $t_i$ to the running
  value.
  \item The second form of \eqref{eq:comb-cocycle} says that a bare vertex
  inserted at depth $M_i$ contributes $t_ie^{-M_i}$ to the source-normalized
  charge and is then transported to the target frame by $e^{M_*}$.
\end{enumerate}

\begin{proposition}[The distributive law is the renormalisation resummation]
\label{prop:distributive-resummation}
Fix two vertices at distinct depths.  Then for every $M_1<M_*$ and $t_1$,
\[
  (\cM_{M_*-M_1},\cA_{t_1})\quad\text{and}\quad
  (\cA_{t_1e^{M_*-M_1}},\cM_{M_*-M_1})
\]
are histories with the same terminal operator but different total charges.  The
difference of their additive charges is the elementary order defect
$t_1\,(e^{M_*-M_1}-1)$.
\end{proposition}

\begin{proof}
Both words evaluate to the affine map $z\mapsto e^{M_*}z+e^{M_*-M_1}t_1$, which
is the content of the pointwise identity $k(z+p)=kz+kp$ of Paper~0,
Subsection~2.3.  Their additive charges are $t_1$ and $t_1e^{M_*-M_1}$, whose
difference is as displayed; the multiplicative charges agree.
\end{proof}

The proposition is the algebraic form of the statement that a large bare vertex
inserted deep in the ultraviolet can be transported to the infrared as a small
effective vertex.  In field-theoretic language the two words are the same
renormalized amplitude with the counterterm absorbed, and the charge difference
is the order defect that the \ACS{} will record as an area.

\subsection{The image in the accumulative commutative space}

Commutativizing forgets the order of the steps and keeps the two total charges
(Paper~0, Definition~8.1 and Proposition~8.5).  For the $p_0\to p$ family this has
an immediate consequence.

\begin{proposition}[The family fans out into a horizontal segment]
\label{prop:acs-fanout}
Let $\{\gamma\}$ be the histories of Proposition~\ref{prop:operator-fiber}.
Their \ACS{} charge chains $C_\gamma$ all terminate on the horizontal line
$M=M_*$, where $M_*$ is fixed by \eqref{eq:fiber-invariants}, and they satisfy
the single scalar constraint
\begin{equation}\label{eq:fanout-constraint}
  \hat\xi(p):=\int_{C_\gamma}e^{-M}\mathrm dA
  =\xi_*\,e^{-M_*}+a(q)
\end{equation}
of Paper~0, Proposition~8.2.  Along that line the total additive charge
$A_\gamma$ is not determined by $p$: a history that inserts its vertices early
needs $A_\gamma\to$ small, and one that inserts them late needs
$A_\gamma\asymp e^{M_*}\hat\xi(p)$.
\end{proposition}

\begin{proof}
The multiplicative charge is determined by the endpoint by
Theorem~8.6(3).  Proposition~8.2 identifies the weighted integral with the
source-normalized terminal translation, which is fixed by \eqref{eq:fiber-invariants}
because both ends are.  The two extremes are the two words of
Proposition~\ref{prop:distributive-resummation}: inserting at $M=0$ costs
$A=t_1$ and inserting at $M=M_*$ costs $A=t_1e^{M_*}$.
\end{proof}

\begin{corollary}[Two readings of the same divergence]
\label{cor:bare-divergence}
The quantity $A_\gamma$ --- the total additive charge, which the \ACS{} keeps as
its abscissa --- is not a function of the physical endpoint, and it grows like
$e^{M_*}$ for histories whose insertions are pushed to the infrared end.  The
weighted quantity $\hat\xi(p)$ is a function of the endpoint and stays finite.
\end{corollary}

\begin{warning}[What this does and does not explain]
\label{warn:divergence}
Corollary~\ref{cor:bare-divergence} says that the unweighted sum of insertions
over a scale range diverges like $e^{M}$ while its $e^{-M}$-weighted version does
not.  That is exactly the algebraic shape of the statement that a bare quantity
regularized at a short distance diverges while a renormalized quantity does not,
and the exponential weight is the usual scale factor.  It is \emph{not} a proof
of renormalizability of any field theory, and no loop integral is computed here.
The correspondence is between one exact \AEG{} identity and one structural
feature of the perturbative expansion.
\end{warning}

When two histories of the family also share their additive charge, they close up
in $\ACS$, and the weighted torsion--Stokes theorem of Paper~0, Theorem~8.9,
applies: the relative defect is
$\tau(\gamma,\delta)=\iint_{\Sigma_{\gamma\delta}}e^{M_*-M}\mathrm dA\wedge
\mathrm dM$.  For two such histories to \emph{also} reach the same $p$, that
torsion must vanish.  In the physical reading this is the statement that
cancellations between early and late insertions --- the finite remainders of
renormalization and the Ward-type identities --- are exactly the vanishing of a
signed weighted area.  The elementary case is the kernel word of
Theorem~8.6: the two histories $\gamma_1=(\cM_q,\cA_p)$ and
$\gamma_2=(\cA_{e^{-q}p},\cM_q)$ reach the same state and differ in additive
charge by $p(1-e^{-q})$, which is exactly the weighted area they enclose
(Lemma~\ref{lem:kernel-word}).

\subsection{The image in the contact model: Reeb heights}

Lift each history horizontally to $\Cont$: a curve $\tilde\gamma$ with
$\alpha(\dot{\tilde\gamma})=0$.  By Paper~0, Theorem~3.9 and the coordinate
form of the horizontal fields, this is the ordinary differential equation
$\mathrm da=\mu\,\mathrm du+\lambda a\,\mathrm dv$, and by
Proposition~9.8(2) its image under $Q$ is the arithmetic motion itself.

\begin{proposition}[Same state, different Reeb height]
\label{prop:reeb-heights}
Let $\tilde\gamma_1,\tilde\gamma_2$ be horizontal lifts of two histories of the
$p_0\to p$ family, starting at the same point of $\Cont$.  Then
$Q\circ\tilde\gamma_1$ and $Q\circ\tilde\gamma_2$ end at the same $p\in\HH^2$,
while
\[
  u(\tilde\gamma_1)-u(\tilde\gamma_2)=\frac{1}{\mu}\bigl(A_{\gamma_1}-A_{\gamma_2}\bigr).
\]
By Theorem~9.9 this height difference equals
$-\iint_{\Omega_{12}}\mathrm da\wedge\mathrm dM=-\mu\lambda\,\mathrm{Area}_g(\Omega_{12})$,
the hyperbolic area enclosed by the two projected paths in $\HH^2$.
\end{proposition}

\begin{proof}
Horizontal lifts of paths with the same endpoints in $\HH^2$ end on the same
fiber of $Q$ because $Q$ is the projection; the height difference is the
statement of Theorem~9.9(2), and the additive charge is $A=\mu u$ by the natural
charge coordinates of Appendix~D.4.
\end{proof}

\begin{figure}[t]
\centering
\begin{tikzpicture}[x=1cm,y=1cm,>=Latex,every node/.style={font=\scriptsize}]
  % ---- panel (a)
  \begin{scope}[shift={(0,0)}]
    \fill[blue!2] (-0.4,-0.4) rectangle (4.0,3.0);
    \draw[black!35] (-0.4,0)--(4.0,0);
    \fill (0.5,2.45) circle (1.7pt) node[above right] {$p_0$ (ultraviolet)};
    \fill (3.3,0.9) circle (1.7pt) node[below right] {$p$};
    \draw[very thick,red!65!black]   (0.5,2.45) to[out=-40,in=170] (3.3,0.9);
    \draw[very thick,blue!65!black]  (0.5,2.45) to[out=-85,in=215] (3.3,0.9);
    \draw[very thick,green!45!black] (0.5,2.45) to[out=-62,in=250] (3.3,0.9);
    \node[font=\small] at (1.8,-0.95) {(a) $\HH^2$};
  \end{scope}
  % ---- panel (b)
  \begin{scope}[shift={(5.3,0)}]
    \fill[blue!2] (-0.5,-0.4) rectangle (3.6,3.0);
    \draw[->,black!45] (-0.4,0.15)--(3.5,0.15) node[right] {$A$};
    \draw[->,black!45] (0,-0.3)--(0,2.85) node[above] {$M$};
    \draw[black!55,dashed] (0,2.2)--(3.3,2.2) node[right] {$M_*$};
    \draw[very thick,red!65!black]   (0,0.15)--(0.5,0.15)--(0.5,2.2);
    \draw[very thick,blue!65!black]  (0,0.15)--(1.4,0.15)--(1.4,2.2);
    \draw[very thick,green!45!black] (0,0.15)--(2.5,0.15)--(2.5,2.2);
    \fill (0.5,2.2) circle (1.6pt);
    \fill (1.4,2.2) circle (1.6pt);
    \fill (2.5,2.2) circle (1.6pt);
    \node[font=\small] at (1.55,-0.95) {(b) $\ACS$};
  \end{scope}
  % ---- panel (c)
  \begin{scope}[shift={(9.9,0)}]
    \fill[blue!2] (-0.7,-0.4) rectangle (2.6,3.0);
    \draw[black!55] (0.8,0.0)--(0.8,2.5) node[above,font=\tiny] {$Q^{-1}(p)$};
    \draw[->,black!45] (2.0,0.0)--(2.0,2.7) node[right,font=\tiny] {$u$};
    \fill[red!65!black]   (0.8,1.2) circle (1.8pt);
    \fill[blue!65!black]  (0.8,1.6) circle (1.8pt);
    \fill[green!45!black] (0.8,2.1) circle (1.8pt);
    \draw[{Latex[length=1.4mm]}-{Latex[length=1.4mm]},black!55]
      (0.28,1.2)--(0.28,2.1) node[midway,left,font=\tiny] {$\Delta A$};
    \node[font=\small] at (0.95,-0.95) {(c) Reeb fiber};
  \end{scope}
\end{tikzpicture}
\caption[Three images of one family of histories]%
{Three images of one family of histories.  (a) In $\HH^2$ the histories have the
same two endpoints and therefore the same composite operator.  (b) In $\ACS$ they
share the multiplicative charge $M_*$ and differ in additive charge $A_\gamma$,
subject to the weighted constraint \eqref{eq:fanout-constraint}.  (c) In $\Cont$
they terminate on one Reeb fiber over $p$ at distinct heights
$u_\gamma=A_\gamma/\mu$; the height spread equals the hyperbolic area enclosed by
the projected paths (Theorem~\ref{thm:dual-tearing}).}
\label{fig:fanout}
\end{figure}

\begin{remark}[Action and phase]
\label{rem:prequantum}
The Reeb coordinate $u$ is the coordinate along which the contact form evaluates
to $1$: $\alpha(R)=1$ with $R=-\tfrac1\mu\pt_u$ (Appendix~D.4 of Paper~0).  In
the geometric quantization of a contact manifold, the Reeb direction is the
circle fiber of the prequantum bundle and its coordinate is the phase of the
wave function \cite{Geiges2008,BoothbyWang1958}.  Proposition~\ref{prop:reeb-heights}
therefore says that histories reaching the same state with different Reeb
heights differ by a phase.  Summing over the fiber with weights
$e^{\,\mathrm i\,A_\gamma}$ is then the natural candidate for a path integral in
this picture.  Nothing in this note depends on that candidate, which is recorded
as a target rather than as a result.
\end{remark}

\subsection{Radix shells}

If the scaling steps are quantized on a $k$-ary grid and each shell either emits
or does not, the word becomes a Horner expansion and the terminal translation is
\begin{equation}\label{eq:horner}
  \xi=\sum_{i=1}^{n}d_i\,k^{\,n-i},
  \qquad d_i\in\{0,1,\dots,k-1\},
\end{equation}
with the counting statements of Paper~IV, Proposition~10.1 and Theorem~10.3.  In
the physical reading $d_i\ne0$ at shell $i$ records a virtual process at the
Wilsonian momentum shell $k^{n-i}$, and the terminal value is built up shell by
shell from the ultraviolet end.  Once signed and continuous increments are
allowed, the number of histories sharing one terminal operator is infinite,
which is the content of Proposition~\ref{prop:operator-fiber}.

\subsection{Panorama}

\begin{center}\small
\renewcommand{\arraystretch}{1.25}
\begin{tabular}{@{}p{0.19\linewidth}p{0.24\linewidth}p{0.24\linewidth}p{0.23\linewidth}@{}}
\toprule
carrier & \AEG{} object & Feynman reading & invariant that survives \\
\midrule
$\E_0$ & broken-line family from $q=\infty$ to $p$; one fiber
$\rho^{-1}(g_{p_0\to p})$
& time- and scale-ordered comb diagrams from the asymptotic state to the dressed
state
& endpoint assignment $a(p)$ and total scale $e^{M_*}$ fixed; the order of
insertions is retained \\
$\ACS$ & all chains end on $M=M_*$; abscissa $A_\gamma=\sum t_i$ free subject to
\eqref{eq:fanout-constraint}
& bare versus renormalized charge; ultraviolet sensitivity of the unweighted sum
& $\hat\xi(p)=\int_{C}e^{-M}\mathrm dA$; and
$\tau(\gamma,\delta)=0$ when both charge and endpoint agree \\
$\Cont$ & Legendrian lift family; endpoints on distinct heights of
$Q^{-1}(p)$
& different histories for one physical state, carrying different phases
& Reeb drift $A_{\gamma_1}-A_{\gamma_2}
=-\mu\lambda\,\mathrm{Area}_g(\Omega_{12})$ \\
\bottomrule
\end{tabular}
\end{center}

Section~\ref{sec:legendre} isolates the last row as a theorem about pairs of
histories and then re-reads the first row as its conjugate.
